% =============================================================================================
%  La structure de l'algorithme de vpp_acceleration, simplement. Note pédagogique (~7 pages),
%  état du code du 8e run (29/10/2026). Pour comprendre la chaîne en entier avant d'entrer dans
%  selection_points_vpp.tex (le choix des calculs) et algorithme_vpp_cible.tex (la référence).
%  Compilation : tectonic -X compile structure_vpp_acceleration.tex --outdir build
% =============================================================================================
\documentclass[11pt,a4paper]{article}
\usepackage{fontspec}
\usepackage[french]{babel}
\usepackage{csquotes}
\usepackage[a4paper,margin=2.0cm,top=1.6cm,bottom=1.8cm]{geometry}
\usepackage{microtype}
\usepackage{amsmath,amssymb,mathtools,bm}
\usepackage{xcolor}
\usepackage{booktabs,tabularx,array,xltabular}
\usepackage{enumitem}
\usepackage{caption}
\usepackage[most]{tcolorbox}
\usepackage{titlesec}
\usepackage{tikz}
\usetikzlibrary{positioning,arrows.meta,calc,fit,backgrounds}
\usepackage[hidelinks,unicode]{hyperref}

\hypersetup{colorlinks=true,linkcolor=blue!50!black,citecolor=green!40!black,
            urlcolor=blue!60!black,pdftitle={La structure de vpp_acceleration}}
\setlength{\parindent}{0pt}
\setlength{\parskip}{0.4em}
\setlength{\emergencystretch}{3em}   % long \texttt paths cannot be hyphenated
\setlist{itemsep=0.1em,topsep=0.2em,leftmargin=1.3em}
\renewcommand{\arraystretch}{1.12}
\titleformat{\section}{\normalfont\large\bfseries\color{blue!40!black}}{\thesection}{0.6em}{}
\titlespacing*{\section}{0pt}{0.8em}{0.25em}
\captionsetup{font=small,labelfont=bf,margin=0.3cm}

% encadrés du document Doc-GP-mfego (preambule.tex)
\tcbset{enhanced,breakable,boxrule=0.6pt,arc=1.5mm,left=2mm,right=2mm,top=0.8mm,bottom=0.8mm,
        fonttitle=\bfseries\small\hypersetup{citecolor=yellow!60!white,linkcolor=yellow!60!white},
        before skip=0.6em,after skip=0.6em}
\newtcolorbox{theorie}[1][]{colback=blue!3,colframe=blue!45!black,
  title={Théorie\ifx&#1&\else~— #1\fi}}
\newtcolorbox{danslecode}[1][]{colback=green!3,colframe=green!35!black,
  title={Dans le code\ifx&#1&\else~— #1\fi}}
\newtcolorbox{subtilite}[1][]{colback=orange!4,colframe=orange!70!black,
  title={Subtilité\ifx&#1&\else~— #1\fi}}
\newtcolorbox{aretenir}[1][]{colback=violet!4,colframe=violet!60!black,
  title={À retenir\ifx&#1&\else~— #1\fi}}

\newcommand{\vx}{\bm{x}}
\newcommand{\vz}{\bm{z}}
\newcommand{\va}{\bm{a}}
\newcommand{\vc}{\bm{c}}
\newcommand{\vs}{\bm{s}}
\newcommand{\vh}{\bm{h}}
\newcommand{\tp}{^{\mathsf{T}}}
\newcommand{\py}[1]{\texttt{#1}}
\DeclareMathOperator{\Var}{Var}
\DeclareMathOperator*{\argmax}{arg\,max}

\begin{document}

{\centering
{\Large\bfseries\color{blue!40!black} La structure de \py{vpp\_acceleration}, simplement\par}
\vspace{0.15em}
{\large Chercher la section du C-foil qui fait aller le bateau le plus vite\par}
\vspace{0.2em}
{\small Note pédagogique · état du code du 8\textsuperscript{e} run (29/10/2026)\par}}

\vspace{0.4em}
\textbf{Ce que cette note explique.} Toute la chaîne, des codes d'hydrodynamique jusqu'à la meilleure
section, en quatre briques : la section (§\,2), le modèle des efforts (§\,3), la VPP (§\,4) et la boucle qui
les relie (§\,5). Chaque brique dit \emph{ce qu'elle fait}, \emph{pourquoi}, et \emph{où elle est dans le
code}. Les détails du choix des nouveaux calculs sont dans \py{selection\_points\_vpp.tex} ; la description
complète, dans \py{algorithme\_vpp\_cible.tex}.

% ---------------------------------------------------------------------------------------------
\section{Le problème en une page}
% ---------------------------------------------------------------------------------------------
On garde le planform du C-foil de l'Intens SY et l'on cherche la \textbf{section 2D} qui maximise la vitesse
du bateau sur ses conditions d'usage, avec une section au moins aussi raide en flexion que l'actuelle.

La troisième colonne donne la notation de la note de M. Le Maître (version 01, 29/10/2026).

\begin{center}\small
\begin{tabularx}{\linewidth}{@{}l>{\raggedright\arraybackslash}Xl@{}}
\toprule
Symbole & Sens & Note Le Maître \\
\midrule
$\vz\in\mathbb R^{10}$ & la section : ses 10 premiers scores PCA (§\,2) & $u_f$ \\
$\va=(\text{cant},\text{rake},\text{yaw})$ & l'attitude effective du foil, en degrés : \emph{un état de fonctionnement}, pas de la géométrie & $z=(\theta_1,\theta_2,\theta_3)$ \\
$\vx=(\vz,\va)\in\mathbb R^{13}$ & l'entrée des modèles d'efforts & $x=(u_f,z)$ \\
$\vh(\vx)=(C_1,\dots,C_7)$ & le torseur $F_x,F_y,F_z,M_x,M_y,M_z$ et le $C_{p,\min}$ (cavitation), en coefficients, donnés par les codes & $\mathsf F$ \\
$\vc=(\text{TWA},\text{TWS})$, $w_c$ & une condition de navigation et sa probabilité ($\sum_c w_c=1$) & $v_w^{(k)}$, $\omega_k$ \\
$\vs$ & les 7 inconnues de la VPP : vitesse, dérive, gîte, assiette, déplacement (l'état), safran, \enquote{flat} (les réglages) ; l'attitude du foil en découle, $\va=A(\vs)$ & $y=(u_r,a)$ \\
$V^\ast(\vz,\vc)$ & la vitesse d'équilibre du bateau avec la section $\vz$, à la condition $\vc$ & $V^\star_k(u_f;\mathsf F)$ \\
$J(\vz)$ & l'objectif : la vitesse moyenne sur le profil d'usage & $J_{\text{moy}}$ \\
$d_k=\partial V^\ast/\partial C_k$ & la sensibilité de la vitesse optimisée à la sortie $k$ (§\,4) & $b_N=B\tp\bar\lambda_N$ \\
\bottomrule
\end{tabularx}
\end{center}

\begin{theorie}[le problème, écrit une fois]
\begin{equation}\label{eq:vstar}
V^\ast(\vz,\vc)=\max_{\vs}\;V(\vs)\quad\text{tel que}\quad
\bm R\big(\vs;\;\vh(\vz,A(\vs)),\;\vc\big)=\bm 0,
\end{equation}
où $\bm R$ regroupe les 6 équations d'équilibre (forces et moments) du bateau entier : voiles, coque, safran
et foil. Puis
\begin{equation}\label{eq:objectif}
\max_{\vz}\;J(\vz)=\sum_c w_c\,V^\ast(\vz,\vc)\qquad\text{tel que}\quad I_{xx}(\vz)\ge I_{xx}(\text{section actuelle}).
\end{equation}
\end{theorie}

\textbf{Qu'est-ce qui coûte ?} Les \textbf{codes} qui calculent $\vh$ : NPLLT × NeuralFoil (environ 0,5 s
par calcul) et surtout AVL × XFOIL itéré (environ 17 s ; chiffres du 8\textsuperscript{e} run). Or la VPP
appelle le modèle du foil à \emph{chaque itération} de son solveur, soit des centaines de fois par solve, pour
chaque condition et chaque section essayée. Brancher les codes directement dans la VPP coûterait des centaines
de calculs AVL par vitesse. D'où l'idée de la chaîne : \textbf{remplacer les codes par un modèle statistique de
$\vh$} (un solve prend alors quelques secondes), et ne lancer les codes que là où ils servent la décision.

\begin{subtilite}[pourquoi ce n'est pas une simple composition $f=g\circ h$]
On peut écrire $J=g\circ\vh$ comme Astudillo \& Frazier~\cite{AF19} : $\vh$ est cher, $g$ (la VPP) est connue
et rapide. Mais chez eux, $\vh$ est évaluée au même point que l'objectif. Ici, la VPP a besoin de $\vh$ en des
attitudes $\va=A(\vs)$ qui sont \emph{elles-mêmes des inconnues de l'équilibre}. Le modèle de $\vh$ doit donc
être bon sur une surface de l'espace $(\vz,\va)$ que l'on ne connaît qu'en résolvant la VPP. C'est la raison
d'être du critère ciblé (§\,5, étape 3).
\end{subtilite}

% ---------------------------------------------------------------------------------------------
\section{Brique 1 : la section}
% ---------------------------------------------------------------------------------------------
\begin{enumerate}
  \item \textbf{16 poids de Kulfan} (8 par face), bord d'attaque arrondi, épaisseur au bord de fuite fixe
        (0,00197 de la corde). C'est la paramétrisation de NeuralFoil.
  \item \textbf{Une famille admissible} : les poids dans une boîte centrée sur la section actuelle, et trois
        contraintes physiques. Les deux faces ne se croisent pas ; $I_{xx}\ge I_{xx}$ actuel ; NeuralFoil
        est confiant (au moins 0,90 à 0°, 2°, 4° et 6°).
  \item \textbf{Une PCA sur les coordonnées} de 4\,000 sections tirées dans la famille. On garde $k=10$
        modes, qui portent 99,87 \% de la variance géométrique. La section est alors $\vz\in\mathbb R^{10}$, et
        le décodage $\vz\mapsto$ poids est exact.
\end{enumerate}

\begin{danslecode}
\py{src/section\_pca.py} : \py{KulfanBoxFamily} (la boîte et les contraintes), \py{SectionPCA} (\py{fit},
\py{encode}, \py{decode}). La taille de la boîte est l'option \py{--box-scale} de
\py{scripts/run\_real\_vpp.py} : 0,5 aux runs 2 à 6, 0,75 au 7\textsuperscript{e}, 3 au 8\textsuperscript{e}.
À l'échelle 3, la boîte ne limite plus : c'est $I_{xx}$ qui retient la section.
\end{danslecode}

% ---------------------------------------------------------------------------------------------
\section{Brique 2 : le modèle des efforts (7 co-krigeages)}
% ---------------------------------------------------------------------------------------------
\textbf{Deux niveaux de fidélité}, calculés aux mêmes points $\vx$ (plan imbriqué) :
\begin{itemize}
  \item niveau 1 (BF) : NPLLT × NeuralFoil, coût relatif 1 ;
  \item niveau 2 (HF) : AVL × XFOIL itéré, coût relatif 6. C'est la référence que la VPP utilise.
\end{itemize}

\textbf{Un co-krigeage AR(1) par sortie} (Kennedy \& O'Hagan~\cite{KOH00}, sous la forme récursive de Le
Gratiet~\cite{LG}) :
\[
Z_2(\vx)=\rho\,Z_1(\vx)+\delta(\vx),
\]
avec un noyau Matérn 5/2 et des hyperparamètres estimés par \textbf{MAP} : prior Inverse-Gamma(3, 2) sur les
longueurs de corrélation, Log-Normale sur la variance et le bruit, $\mathcal N(1,2^2)$ sur $\rho$. Il y a 7
co-krigeages indépendants, un par composante de $\vh$, qui partagent le même plan.

\begin{aretenir}[ce que les GP apprennent, et ce qu'ils n'apprennent pas]
\begin{itemize}
  \item Les GP apprennent $\vh(\vz,\va)$ \textbf{directement}, à partir des sorties des codes, en des points
        $(\vz,\va)$ que nous choisissons. C'est une \textbf{régression ordinaire}, sortie par sortie.
  \item \textbf{Aucun GP n'est appris à travers l'équation d'équilibre.} On n'observe jamais $V^\ast$ pour en
        déduire $\vh$. L'équilibre ne sert qu'\emph{après} l'apprentissage : pour prédire la vitesse (§\,4) et
        pour choisir les prochains calculs (§\,5).
\end{itemize}
\end{aretenir}

\begin{danslecode}
\py{src/vpp\_targeted.py} : \py{TorsorCoKriging} (les 7 \py{CoKriging} de mfego, ajustement, prédiction,
conditionnement). Les codes : \py{src/foil\_codes.py} (sur bdfoil\_v1), lancés en parallèle par
\py{CodeRunner} (\py{src/vpp\_optimize.py}).
\end{danslecode}

% ---------------------------------------------------------------------------------------------
\section{Brique 3 : la VPP et le profil d'usage}
% ---------------------------------------------------------------------------------------------
\begin{itemize}
  \item \textbf{La VPP} est celle de MDS (Intens SY, foil \enquote{CFoilFSI}). Elle maximise la vitesse sous les
        6 équations d'équilibre, avec 7 inconnues (vitesse, dérive, gîte, assiette, déplacement, safran,
        flat). Nous remplaçons seulement la surface de réponse du foil par \textbf{la moyenne} des 7
        co-krigeages. Une correction de déformation du foil (FSI) s'y ajoute, mise à l'échelle par la raideur
        de la section.
  \item \textbf{Le profil d'usage} donne les conditions $\vc$ et leurs poids $w_c$. Au
        8\textsuperscript{e} run, c'est \py{profile\_v2.json} (dossier \py{profiles/}), avec 49 conditions, deux bosses à 40° et 80° de TWA, TWS de 5 à
        17 nd, et des points sur les bords du domaine. Il se construit avec l'éditeur HTML du dossier
        \py{Bimodal-Sailing-Discretizer}.
\end{itemize}

\begin{theorie}[l'incertitude sur la vitesse, au premier ordre]
À un état d'équilibre, une petite erreur $\delta C_k$ change la vitesse de $d_k\,\delta C_k$, avec
$d_k=\partial V^\ast/\partial C_k$ calculé par différences finies avec la VPP. Les 7 GP étant indépendants,
\[
\Var V^\ast\;\approx\;\sum_{k} d_k^2\,s_k^2(\vx),
\]
avec $s_k^2$ la variance du co-krigeage $k$. Une erreur de traînée coûte beaucoup en vitesse ; une erreur
d'effort latéral presque rien, car la VPP rattrape l'équilibre avec une dérive un peu différente.
\end{theorie}

\begin{subtilite}[deux limites à connaître]
\begin{itemize}
  \item \textbf{On met la moyenne dans la VPP}, alors que $\mathbb E[V^\ast(\vh)]\ne V^\ast(\mathbb E[\vh])$ si
        la VPP est non linéaire. L'écart est du second ordre en $s_k$ : il est négligeable là où le modèle est
        précis. Au 8\textsuperscript{e} run, l'erreur de vitesse contre AVL × XFOIL est de 0,0006 nd à la
        meilleure section.
  \item \textbf{La VPP de MDS bute à 24 nd} : les surfaces de réponse de la coque ne vont pas plus loin. Au
        travers par vent fort, plusieurs conditions sont en butée et ne départagent pas les sections.
\end{itemize}
\end{subtilite}

\begin{danslecode}[comment $d_k$ est calculé, honnêtement]
\begin{itemize}
  \item \textbf{Une différence finie avant, à un seul pas.} On ajoute une constante à la sortie $k$ du modèle,
        partout (environ 3 \% de sa taille : 5e-4 pour $F_x$, 1e-2 pour $F_y$ et $F_z$, 2e-2 pour $M_x$, 3e-3
        pour $M_y$ et $M_z$), on refait la VPP complète en partant de la base, et
        $d_k=(V_{\text{décalé}}-V_{\text{base}})/\text{pas}$. Toutes les inconnues sont réoptimisées : c'est la
        sensibilité de la vitesse \emph{optimisée}.
  \item \textbf{12 solves par condition} : pour chaque $k$, la base est refaite dans le même processus que le
        solve décalé, car le solveur MDS dépend du chemin.
  \item \textbf{Seulement à la titulaire}, une mise à jour sur trois, puis la médiane des 3 derniers calculs. Un
        solve en échec garde l'ancienne valeur ; une condition en butée à 24 nd a $d=0$.
  \item \textbf{Jamais vérifié} contre un autre pas, une différence centrée ou les multiplicateurs (§\,9).
\end{itemize}
Le code : la méthode \py{speed\_sensitivities} de \py{IntensSYVPP} (\py{src/vpp\_coupling.py}) et la
méthode \py{sensitivities} de \py{VPPSectionOptimizer} (\py{src/vpp\_optimize.py}). Le solve est
\py{IntensSYVPP.solve} ; le profil, \py{OperatingProfile} (\py{src/operating\_profile.py}).
\end{danslecode}

% ---------------------------------------------------------------------------------------------
\section{Brique 4 : la boucle}
% ---------------------------------------------------------------------------------------------
\begin{figure}[h]
\centering
\begin{tikzpicture}[
  node distance=0.55cm and 1.15cm,
  box/.style={draw,rounded corners=2pt,align=center,font=\small,minimum height=1.0cm,text width=3.15cm,fill=#1},
  box/.default=white,
  arr/.style={-{Stealth[length=2.2mm]},thick}]
\node[box=gray!10] (codes) {\textbf{Codes}\\NPLLT × NF (BF)\\AVL × XFOIL (HF)};
\node[box=blue!6, right=of codes] (gp) {\textbf{7 co-krigeages}\\$\vh(\vz,\va)$ : moyenne $m_k$, variance $s_k^2$};
\node[box=green!6, right=of gp] (vpp) {\textbf{VPP} (moyenne $m_k$)\\vitesse $V^\ast$, états $\vs$, sensibilités $d_k$};
\node[box=orange!8, below=1.1cm of vpp] (crib) {\textbf{1–2. Criblage}\\256 sections au 1\textsuperscript{er} ordre, 3 à la VPP : le titulaire};
\node[box=violet!6, left=of crib] (crit) {\textbf{3–4. Critère ALC ciblé}\\nuages autour des états, choix (point, niveau) par gain / coût};
\node[box=gray!10, left=of crit] (lot) {\textbf{5. Lot de 16 calculs}\\\enquote{kriging believer}, en parallèle};
\draw[arr] (codes) -- node[above,font=\scriptsize]{données} (gp);
\draw[arr] (gp) -- (vpp);
\draw[arr] (vpp) -- (crib);
\draw[arr] (crib) -- node[above,font=\scriptsize]{états} (crit);
\draw[arr] (crit) -- (lot);
\draw[arr] (lot) -- (codes);
\end{tikzpicture}
\caption{Une mise à jour de la boucle. Les codes ne sont lancés qu'à l'étape 5 ; tout le reste utilise le
modèle et la VPP.}
\label{fig:boucle}
\end{figure}

\textbf{Étape 0 — le plan de départ.} 256 points BF : des sections de la famille, avec une attitude tirée
par hypercube latin. Puis 64 points HF choisis par maximin parmi eux (plan imbriqué). On ajuste les 7
co-krigeages. La recherche part de la section actuelle (ou de la section admissible la plus proche).

Ensuite, à chaque mise à jour :

\textbf{Étape 1 — cribler des sections.} On tire 256 sections : la moitié dans une \textbf{région de
confiance} autour du titulaire (la meilleure section du moment), l'autre dans toute la famille. Pour chacune,
on prédit au premier ordre le gain de vitesse $\mu$ et son incertitude $\sigma$, sans résoudre la VPP. On
garde les 3 meilleurs scores $\mu+\sigma$ : une section incertaine mais prometteuse a sa chance (comme
UCB~\cite{Srinivas10}).

\textbf{Étape 2 — juger avec la VPP complète.} Les 3 sections retenues et le titulaire passent à la VPP, sur
toutes les conditions. La meilleure vitesse moyenne pondérée devient le titulaire. La région de confiance
grandit (×1,5) si le titulaire a changé, et rétrécit (×0,7) sinon (TuRBO~\cite{TuRBO19}).

\textbf{Étape 3 — où le modèle doit être bon.} Autour de chaque état d'équilibre de ces 4 sections, un nuage
de 16 points : l'attitude bouge de $\pm1°$ en cant, $\pm0{,}3°$ en rake et $\pm1°$ en yaw, et la section aussi
sur un point sur deux. Chaque point pèse $w_c\,d_k^2$ : la probabilité de sa condition, fois l'effet de
l'erreur sur la vitesse.

\textbf{Étape 4 — choisir les calculs (le critère ALC ciblé).} Pour chaque candidat $\vx_m$ et chaque niveau
$l$, on sait \emph{avant de lancer le code} de combien la variance pondérée des nuages baissera : la variance
d'un GP ne dépend pas des valeurs observées. On prend le meilleur rapport
\[
(\vx^\star,l^\star)=\argmax_{m,\,l}\;\frac{\text{baisse de }\sum\text{(poids)}\times\text{variance sur les nuages}}{\text{coût du calcul}}.
\]
C'est l'ALC de Cohn~\cite{Cohn96} et de Gramacy \& Lee~\cite{GL09}, avec la mesure d'intégration placée
autour des équilibres au lieu de tout le domaine (Le Gratiet \& Cannamela~\cite{LGC15}).

\textbf{Étape 5 — 16 calculs à la fois.} On choisit 16 couples (point, niveau) l'un après l'autre. Après
chaque choix, on fait comme si le calcul était fait (\enquote{kriging believer}~\cite{GLC10}) : le choix
suivant voit exactement la variance déjà retirée, et le lot s'étale de lui-même. Les 16 calculs tournent en
parallèle, puis les vraies valeurs entrent dans le modèle.

\textbf{Étape 6 — s'arrêter ou repartir.} Si les 10 dernières mises à jour ont gagné moins de 0,01 nd en
tout, la région a convergé : on garde son optimum local et l'on repart (au plus 2 fois) d'une section éloignée
de tous les optima, en gardant toutes les données. La boucle s'arrête aussi à la limite de temps.

\textbf{Étape 7 — vérifier.} La meilleure section, la section de départ, les autres optima locaux et les
meilleures sections des runs précédents sont recalculés avec les codes aux états VPP. On compare la vitesse
du modèle à celle des codes, puis l'on analyse les sections sur le profil d'usage. Les coordonnées des
sections (format Selig) vont dans \py{results/sections/}, les figures dans \py{figures/}.

\begin{danslecode}
\py{scripts/run\_real\_vpp.py} enchaîne tout (étapes 0 à 7) et écrit un dossier \py{runs/<date>\_<nom>/}.
La boucle est \py{VPPSectionOptimizer} (\py{src/vpp\_optimize.py}) ; le critère, \py{choose}
(\py{src/vpp\_targeted.py}) et les gains exacts de \py{src/sequential\_design.py} ; les figures,
\py{scripts/real\_run\_figures.py}.
\end{danslecode}

% ---------------------------------------------------------------------------------------------
\section{Ce qui ne sert qu'à aller vite}
% ---------------------------------------------------------------------------------------------
Ces couches ne changent pas la logique. Sans elles, l'algorithme est le même, en plus lent :
\begin{itemize}
  \item les lots de 16 et les processus parallèles (28 pour les codes, 28 pour la VPP sur toute la machine) ;
  \item la mise à jour du facteur de Cholesky au lieu de tout refactoriser, pendant un lot ;
  \item les hyperparamètres réestimés seulement quand le plan a grossi de 25 \%, et plus du tout au-delà de
        2\,500 points ;
  \item les sensibilités $d_k$ recalculées une mise à jour sur trois (médiane des 3 dernières) ;
  \item la VPP avec la moyenne seule (sans variance) quand on n'a besoin que de la vitesse.
\end{itemize}

% ---------------------------------------------------------------------------------------------
\section{Les hypothèses essentielles}
% ---------------------------------------------------------------------------------------------
{\small
\begin{xltabular}{\linewidth}{@{}>{\raggedright\arraybackslash}p{4.6cm}>{\raggedright\arraybackslash}X>{\raggedright\arraybackslash}X@{}}
\toprule
Hypothèse & Pourquoi & Ce qu'on risque \\
\midrule
\endhead
AR(1) à $\rho$ constant, plan imbriqué & le modèle de \cite{LG}, validé sur nos cas & une corrélation BF/HF qui varie dans le domaine \\
7 sorties indépendantes & un calcul donne tout le torseur ; simple & corrélations physiques ignorées ($F_y$ et $M_x$ viennent de la même portance) \\
hyperparamètres par MAP & stable quand il y a peu de données & résultat dépendant du prior au début du run \\
moyenne du modèle dans la VPP & la VPP reste déterministe & biais de second ordre là où le modèle est incertain \\
vitesse linéarisée ($d_k$) & rend le critère calculable & faux si l'erreur déplace beaucoup l'équilibre \\
intégration autour des états de 4 sections & seule zone utile à la décision & le reste du domaine peu appris \\
criblage au premier ordre & trier 256 sections sans 256 VPP & une bonne section mal classée \\
famille sans contrainte de forme (8\textsuperscript{e} run) & ne pas brider l'optimum & une section à l'extrado non convexe \\
\bottomrule
\end{xltabular}}

La liste complète (13 hypothèses) est dans \py{selection\_points\_vpp.tex}, §\,7.

% ---------------------------------------------------------------------------------------------
\section{Où on en est}
% ---------------------------------------------------------------------------------------------
{\small
\begin{tabularx}{\linewidth}{@{}l>{\raggedright\arraybackslash}X r@{}}
\toprule
Run & Ce qui change & Gain sur la section actuelle \\
\midrule
4\textsuperscript{e} (23/10, 8 h) & recherche sur une distribution d'usage, recherche locale & +0,16 nd \\
5\textsuperscript{e} (27/10, 2 h) & convergence, redémarrages, lots par Cholesky : même résultat en 2 h & +0,15 nd \\
6\textsuperscript{e} (27/10, 6 h) & 3 niveaux sur bdfoil\_v1 : le niveau du milieu ne vaut pas son coût & +0,16 nd \\
7\textsuperscript{e} (28/10, 8 h) & profil \py{profile\_v1}, boîte 0,75 : un nouveau bassin & +0,21 nd \\
8\textsuperscript{e} (29/10, 2 h 49) & profil \py{profile\_v2}, boîte 3 (seule $I_{xx}$ retient la section) & +0,26 nd \\
\bottomrule
\end{tabularx}}

Chaque gain est mesuré sur le profil du run, contre AVL × XFOIL. Au 8\textsuperscript{e} run, la section est
moins cambrée (2,2 \% à 55 \% de corde, contre 4,2 \% à 33 \% pour l'actuelle). Elle gagne au travers et perd
un peu au près serré par vent fort. Son extrado n'est plus convexe (creux de 3,6 ‰ de corde à 80 \%) : c'est
à vérifier.

% ---------------------------------------------------------------------------------------------
\section{Ce que nous devrions changer (note de M. Le Maître)}
% ---------------------------------------------------------------------------------------------
La différence principale : \textbf{notre ALC réduit la variance de la vitesse de chaque section}, prise une
par une ; \textbf{son critère réduit le regret de conception}, la vitesse qu'on s'attend à perdre en
choisissant mal. Le regret ne dépend que des écarts de moyenne entre sections candidates et de la variance de
leurs \emph{différences} $S_{ii}+S_{jj}-2S_{ij}$ : une incertitude commune ne vaut rien. Les changements, dans
l'ordre :
\begin{enumerate}
  \item \textbf{Les sensibilités par les multiplicateurs} du VPP, $d=B\tp\bar\lambda$ (sa 5.14) : un solve au
        lieu de 12 par condition, vérifié contre nos différences finies. Cela rend abordable une sensibilité
        $b_i$ \emph{par section candidate}, à son propre équilibre, au lieu de celle de la titulaire pour tout
        le monde.
  \item \textbf{La loi jointe complète} des vitesses candidates, $S_{ij}=\sum_{c,c'}w_cw_{c'}\,b_{ic}\tp
        K_N(x^\star_{ic},x^\star_{jc'})\,b_{jc'}$ (sa 12.6), entre sections \emph{et} entre conditions, au lieu des
        seules variances au criblage.
  \item \textbf{L'arrêt sur le regret}, $\tilde r_N=\mathbb E[\max_i\tilde J_i]-\max_i m_i\le\varepsilon_V$ (en
        nœuds), au lieu de \enquote{moins de 0,01 nd gagné en 10 mises à jour} : on s'arrête quand le choix est
        sûr, pas quand la recherche stagne. Les redémarrages restent : le regret ne voit que les candidates
        connues.
  \item \textbf{Le critère de regret (sa 9.11) à la place de l'ALC}, en version par lot et à deux niveaux,
        comparé à l'ALC sur un cas test.
  \item \textbf{Les hypothèses sur les entrées}, écrites : pas de vitesse $V_f$ (coefficients, pas d'effet de
        Reynolds, sauf la correction FSI) ni d'enfoncement $h_f$ (pas de surface libre dans AVL et XFOIL).
\end{enumerate}

% ---------------------------------------------------------------------------------------------
\section{Questions ouvertes}
% ---------------------------------------------------------------------------------------------
\begin{enumerate}
  \item \textbf{Les hyperparamètres par validation croisée (LOO)} au lieu du MAP, ou en diagnostic : les
        formules rapides existent pour le co-krigeage~\cite{LGC15}, et la LOO en forme fermée est déjà dans
        mfego (\py{CoKriging.leave\_one\_out}). La calibration de la variance compte, puisque le critère, le
        nôtre comme le sien, est fait de variances.
  \item \textbf{Les sorties corrélées} : un modèle multi-sorties au lieu de 7 GP indépendants.
  \item \textbf{L'incertitude à travers l'équilibre} : propager la loi de $\vh$ dans la VPP au lieu de la
        moyenne.
  \item \textbf{La forme} : faut-il une contrainte de convexité ou de régularité dans la famille ?
  \item \textbf{La butée de 24 nd} de la VPP au travers par vent fort.
\end{enumerate}

\begin{thebibliography}{9}\footnotesize\setlength{\itemsep}{0pt}
\bibitem{KOH00} M. C. Kennedy, A. O'Hagan, \emph{Predicting the output from a complex computer code when fast
approximations are available}, Biometrika 87(1), 1–13, 2000.
\bibitem{LG} L. Le Gratiet, \emph{Multi-fidelity Gaussian process regression for computer experiments},
thèse, Université Paris-Diderot, 2013.
\bibitem{LGC15} L. Le Gratiet, C. Cannamela, \emph{Cokriging-based sequential design strategies using fast
cross-validation techniques for multi-fidelity computer codes}, Technometrics 57(3), 418–427, 2015.
\bibitem{Cohn96} D. Cohn, \emph{Neural network exploration using optimal experiment design}, Neural
Networks 9(6), 1996.
\bibitem{GL09} R. B. Gramacy, H. K. H. Lee, \emph{Adaptive design and analysis of supercomputer experiments},
Technometrics 51(2), 130–145, 2009.
\bibitem{AF19} R. Astudillo, P. I. Frazier, \emph{Bayesian optimization of composite functions}, ICML 2019.
\bibitem{GLC10} D. Ginsbourger, R. Le Riche, L. Carraro, \emph{Kriging is well-suited to parallelize
optimization}, in \emph{Computational Intelligence in Expensive Optimization Problems}, Springer, 2010.
\bibitem{TuRBO19} D. Eriksson, M. Pearce, J. Gardner, R. Turner, M. Poloczek, \emph{Scalable global
optimization via local Bayesian optimization}, NeurIPS 2019.
\bibitem{Srinivas10} N. Srinivas, A. Krause, S. Kakade, M. Seeger, \emph{Gaussian process optimization in
the bandit setting: no regret and experimental design}, ICML 2010.
\end{thebibliography}

\end{document}
